\subsection{BST INSERT와 DELETE}
\begin{frame}{BST INSERT의 두 Cursor}
\begin{description}\item[$x$] search cursor: 다음 비교 node\item[$y$] $x$의 parent 후보: 마지막 non-NIL node\end{description}
\mission{새 key가 들어갈 NIL slot을 찾되 parent를 잃지 않는다.}
\end{frame}
\begin{frame}[shrink=2]{TreeInsert Pseudocode}
\scriptsize\begin{algorithmic}[1]
\Procedure{TreeInsert}{$T,z$}
\State $y\gets\NIL$; $x\gets T.root$
\While{$x\ne\NIL$}
\State $y\gets x$
\If{$z.key=x.key$}\State\Return $\Duplicate$\EndIf
\If{$z.key<x.key$}\State $x\gets x.left$\Else\State $x\gets x.right$\EndIf
\EndWhile
\State $z.parent\gets y$
\If{$y=\NIL$}\State $T.root\gets z$
\ElsIf{$z.key<y.key$}\State $y.left\gets z$
\Else\State $y.right\gets z$\EndIf
\EndProcedure
\end{algorithmic}
\end{frame}
\begin{frame}{INSERT 12: 실제 BST의 \texorpdfstring{$x/y$}{x/y}}
\centering\begin{tikzpicture}[x=1.1cm,y=.55cm]
\node[font=\scriptsize,text=AlgoGray]at(-4.4,.2){full BST; orange path만 따라감};
\node[st inactive](n15)at(0,0){15};
\node[st inactive](n6)at(-3,-1){6}; \node[st inactive](n18)at(3,-1){18};
\node[st inactive](n3)at(-4.3,-2){3}; \node[st inactive](n7)at(-1.7,-2){7};
\node[st inactive](n17)at(2.3,-2){17}; \node[st inactive](n20)at(3.8,-2){20};
\node[st inactive](n2)at(-5,-3){2}; \node[st inactive](n4)at(-3.7,-3){4};
\node[st inactive](n13)at(-.5,-3){13};
\node[st inactive](n9)at(-1.2,-4){9}; \node[st inactive](n14)at(.2,-4){14};
\coordinate(n12)at(-.5,-5);
\draw[st edge](n15)--(n6)(n15)--(n18)(n6)--(n3)(n6)--(n7)
 (n3)--(n2)(n3)--(n4)(n7)--(n13)(n13)--(n9)(n13)--(n14)
 (n18)--(n17)(n18)--(n20);
\only<2-|handout:1>{\draw[st active edge](n15)--(n6);}
\only<3-|handout:1>{\draw[st active edge](n6)--(n7);}
\only<4-|handout:1>{\draw[st active edge](n7)--(n13);}
\only<5-|handout:1>{\draw[st active edge](n13)--(n9);}
\only<6|handout:1>{\draw[st active edge](n9)--(n12);\node[st current]at(n12){12};}
\only<1|handout:0>{\node[st current]at(n15){15};\node[callout]at(4.5,-4.3){$x=15,\ y=\NIL$\\$12<15$};}
\only<2|handout:0>{\node[st done]at(n15){15};\node[st current]at(n6){6};\node[callout]at(4.5,-4.3){$x=6,\ y=15$\\$12>6$};}
\only<3|handout:0>{\node[st done]at(n15){15};\node[st done]at(n6){6};\node[st current]at(n7){7};\node[callout]at(4.5,-4.3){$x=7,\ y=6$\\$12>7$};}
\only<4|handout:0>{\node[st done]at(n15){15};\node[st done]at(n6){6};\node[st done]at(n7){7};\node[st current]at(n13){13};\node[callout]at(4.5,-4.3){$x=13,\ y=7$\\$12<13$};}
\only<5|handout:0>{\node[st done]at(n15){15};\node[st done]at(n6){6};\node[st done]at(n7){7};\node[st done]at(n13){13};\node[st current]at(n9){9};\node[callout]at(4.5,-4.3){$x=9,\ y=13$\\$12>9$};}
\only<6|handout:1>{\node[st done]at(n15){15};\node[st done]at(n6){6};\node[st done]at(n7){7};\node[st done]at(n13){13};\node[st done]at(n9){9};\node[callout]at(4.5,-4.3){$x=\NIL,\ y=9$\\9의 right child로 attach};}
\end{tikzpicture}
\end{frame}
\begin{frame}{INSERT Correctness와 비용}
\begin{itemize}
\item search path의 비교가 lower/upper bound를 좁힌다.
\item 마지막 $y$의 올바른 NIL side에 연결하면 모든 ancestor bound 유지
\item path length만큼 비교 → $\Theta(h)$
\end{itemize}
\stcheckpoint{\textbf{INSERT 14를 시도하면?} 기존 node 14에 도달해 $\Duplicate$를 반환하고 tree는 바뀌지 않는다.}
\end{frame}
\begin{frame}{DELETE의 세 Case}
\begin{enumerate}\item child 0개: node 제거\item child 1개: child subtree를 target 위치로 올림\item child 2개: right subtree minimum인 successor로 대체\end{enumerate}
\begin{alertblock}{왜 어려운가?}target 아래의 모든 subtree와 parent/root pointer를 잃지 않아야 한다.\end{alertblock}
\end{frame}
\begin{frame}{Transplant Helper}
\small\begin{algorithmic}[1]
\Procedure{Transplant}{$T,u,v$}
\If{$u.parent=\NIL$}\State $T.root\gets v$
\ElsIf{$u=u.parent.left$}\State $u.parent.left\gets v$
\Else\State $u.parent.right\gets v$\EndIf
\If{$v\ne\NIL$}\State $v.parent\gets u.parent$\EndIf
\EndProcedure
\end{algorithmic}
\sttakeaway{Node $u$의 parent link를 $v$로 바꾸며 root case도 같은 helper가 처리한다.}
\end{frame}
\begin{frame}{DELETE Case 1: Leaf}
\centering\begin{tikzpicture}
\node[st node](p)at(0,0){6};\node[st node](l)at(-1,-1){3};\node[st violation](z)at(1,-1){7};
\draw[st edge](p)--(l);\only<1|handout:0>{\draw[st edge](p)--(z);\node[callout]at(3,-1){target $z=7$};}
\only<2|handout:1>{\node[st nil]at(1,-1){NIL};\node[callout]at(3,-1){parent link를 NIL로};}
\end{tikzpicture}
\end{frame}
\begin{frame}{DELETE Case 2: One Child}
\centering\begin{tikzpicture}
\node[st node](p)at(0,0){15};\node[st violation](z)at(-1.3,-1){6};\node[st done](c)at(-1.3,-2){3};\node[st node](r)at(1.3,-1){18};
\draw[st edge](p)--(r);
\only<1|handout:0>{\draw[st edge](p)--(z)(z)--(c);\node[callout]at(3,-1.4){Transplant(6,3)};}
\only<2|handout:1>{\draw[st active edge](p)--(c);\node[callout]at(3,-1.4){subtree 3이 6의 자리로};}
\end{tikzpicture}
\end{frame}
\begin{frame}{DELETE Case 3: Successor 선택}
\centering\begin{tikzpicture}[x=1.3cm,y=.9cm]
\node[st violation](z)at(0,0){15};\node[st node](l)at(-2,-1){6};\node[st node](r)at(2,-1){20};
\node[st current](y)at(1,-2){17};\node[st node](q)at(3,-2){25};\node[st done](c)at(1.5,-3){18};
\draw[st edge](z)--(l)(z)--(r)(r)--(y)(r)--(q)(y)--(c);
\node[callout,align=center]at(5,-1.5){$y=min(z.right)=17$\\따라서 $y.left=\NIL$};
\end{tikzpicture}
\end{frame}
\begin{frame}{DELETE Case 3: Transplant Sequence}
\small\begin{enumerate}
\item $y=17$, $y.right=18$: \texttt{TRANSPLANT}$(17,18)$ → $20.left=18$
\item $y.right\gets z.right$: $17.right=20$과 parent update
\item \texttt{TRANSPLANT}$(15,17)$
\item $y.left\gets z.left$: $17.left=$ old left subtree와 parent update
\end{enumerate}
\begin{block}{Node identity}key copy가 아니라 successor node 자체를 옮겨 satellite data와 identity를 분명히 한다.\end{block}
\end{frame}
\begin{frame}{DELETE Case 3: Final Tree}
\centering\begin{tikzpicture}[x=1.4cm,y=.9cm]
\node[st current](y)at(0,0){17};\node[st node](l)at(-2,-1){6};\node[st node](r)at(2,-1){20};
\node[st done](c)at(1,-2){18};\node[st node](q)at(3,-2){25};
\draw[st active edge](y)--(l)(y)--(r);\draw[st edge](r)--(c)(r)--(q);
\node[callout]at(5,-1){표시된 tree의 inorder:\\$6,17,18,20,25$};
\end{tikzpicture}
\end{frame}
\begin{frame}{DELETE Correctness Intuition}
\begin{description}
\item[Case 1/2] child subtree 전체가 올라가도 기존 ancestor bounds 안에 있다.
\item[Case 3] successor $y$는 $z$보다 큰 최소 key이고 $y.left=\NIL$이다.
\item[Bounds] $z.left$의 모든 key $<y$; 남은 $z.right$의 모든 key $>y$.
\end{description}
\sttakeaway{삭제 뒤 inorder가 strictly increasing이면 BST order가 보존되었다.}
\end{frame}
\begin{frame}{BST INSERT · DELETE Checkpoint}
\begin{enumerate}\item root 하나뿐인 tree에서 root 삭제는 누가 처리하는가?\item successor가 left child를 가질 수 없는 이유는?\item 두 operation의 시간은?\end{enumerate}
\begin{exampleblock}{답}Transplant의 root branch; right subtree minimum이므로; $\Theta(h)$.\end{exampleblock}
\sttakeaway{BST 연산은 모두 $\Theta(h)$다. 이제 입력 순서가 만든 큰 $h$를 구조적으로 제한한다.}
\end{frame}
